\begin{pastexam}[label={px:ncku-107-cce-08}]{成大電通所 107 學年度 問答題第 4 題（節錄 (a)(b)(c)(e)(f) 小題）}
（問答；14 分，本書收其中 12 分）

\textbf{4. Tree} (14 pts)

Given the following pseudo code whose input includes two trees which are pointed by $s$ and $t$,
respectively, the algorithm checks equality of the two trees.
\begin{minted}[fontsize=\small,escapeinside=@@]{text}
1. bool operator==(const Tree& s, const Tree& t) { // Driver
2.        return Equal(s.root, t.root);
3. }

4. bool Tree <T>::Equal(TreeNode <T> *a, TreeNode <T> *b) { // Workhorse
5.     @\underline{\textbf{Terminal Condition}}@
6.     else return  ( a && b
7.                  && (a->data == b->data)
8.                  && Equal(a->leftChild, b->leftChild)
9.                  && Equal(a->rightChild, b->rightChild));   }
\end{minted}
\begin{enumerate}[label=\alph*.]
\item (\textbf{2 pts}) Please complete the terminal condition in \textbf{Line} 5.
\item (\textbf{2 pts}) What is the purpose of \textbf{Line} 6?
\item (\textbf{3 pts}) Can we obtain a correct answer if we change the sequence of Lines 7 to 9? Why?
\end{enumerate}

The \textbf{satisfiability (SAT)} is an important problem to answer the question if there exists an
assignment to make a boolean expression true.
\begin{minted}[fontsize=\small]{text}
9.  enum Operator {Not, And, Or, True, False};
10. class SatTree; // forward declaration
11. class SatNode {
12.           friend class SatTree;
13.           Operator data;
14.           bool value;
15.           SatNode *leftChild;
16            SatNode *rightChild;     }
\end{minted}
\begin{minted}[fontsize=\small]{text}
17. void SatTree::PostOrderEval() {PostOrderEval(root);}

18. void SatTree::PostOrderEval(SatNode *s) {
19.    if (s) { // not null
          ??????????????????
20.  }
\end{minted}
\begin{enumerate}[label=\alph*.,start=5]
\item (\textbf{2 pts}) What is the correct sequence inside \textbf{Lines} 19-20 according to the following
codes? (A) 甲 乙 丙 (B) 甲 丙 乙 (C) 乙 甲 丙 (D) None of above

甲、\texttt{PostOrderEval(s->leftChild);}\\
乙、
\begin{minted}[fontsize=\footnotesize]{text}
switch(s->data) {
  b.1  case Not:   s->value = !s->rightChild->value; break;
  b.2  case And:   s->value = s->leftChild->value && s->rightChild->value; break;
  b.3  case Or:    s->value = s->leftChild->value || s->rightChild->value;; break;
  b.4  case True:  s->value = true; break; // terminal node
  b.5  case False: s->value = false; // terminal node
}
\end{minted}
丙、\texttt{PostOrderEval(s->rightChild);}
\item (\textbf{3 pts}) Please give the reason for your choice in the last question.
\end{enumerate}
\end{pastexam}

\begin{pastsoln}
這份考卷和成大電機所 107 學年度的「資料結構」是同一份（兩系所共用），完整詳解見歷屆試題~\ref{px:ncku-107-ee-ds-08}。

答案：a. \texttt{if ((!a) \&\& (!b)) return true;}；b. 確認兩個指標都不是空的，恰好一邊是空的就回傳
\texttt{false}，也避免後面對空指標取值；c. 可以，三個條件都沒有副作用，\texttt{\&\&} 的結果和順序無關，
只影響效率；e. (B) 甲 丙 乙；f. 節點的值要用子節點的值算，必須先算左、右子樹再算自己，是後序走訪。
(d) 小題考 C++ 的 \texttt{friend}，本書不收。
\end{pastsoln}
